\documentclass[11pt]{article}
\usepackage[margin=1in]{geometry}
\usepackage[T1]{fontenc}
\usepackage{lmodern,microtype,booktabs,amsmath,graphicx,pdfpages}
\usepackage[colorlinks=true,allcolors=blue!50!black]{hyperref}
\usepackage{xcolor}
\usepackage{caption}
\captionsetup{font=small,labelfont=bf}
\setlength{\parskip}{4pt}
\title{Research and Market Agreement Across Frontier Models:\\An Expanded Prospective Comparison}
\author{Vorhersage experimental series}
\date{September 18, 2026 (America/New\_York)}
\begin{document}
\maketitle
\begin{abstract}
We compare three language models on 15 unresolved prediction-market contracts,
using identical questions and a common set of outside evidence. Each model
produces a question-only forecast and a separate research-informed forecast,
for 90 calls in total. Market prices are committed before research and withheld
from all forecasting inputs. On the 13 contracts with accepted research forecasts
from every model, mean absolute distance from the opening Kalshi midpoint is
8.92 percentage points for Claude Fable 5.1, 11.08 for GPT-6 Astra, and 13.08
for Gemini 3.1 Pro. Equal-weight averaging yields 10.87 points. On the nine
contracts with all six accepted forecasts, research reduces mean distance by
38--45\% across models. An overly broad lexical screen excludes ten otherwise
valid responses; an explicitly post-hoc analysis of the original responses
preserves the aggregate direction. The results support further study of research
augmentation while revealing substantial sensitivity to question composition
and response-screening policy.
\end{abstract}

\section{Objective and experimental design}
Market probabilities offer a useful intermediate target for evaluating forecasts
before events resolve. They incorporate contemporaneous information and permit
rapid comparison without importing realized outcomes into the exercise. This
study asks whether supplying a fixed body of nonmarket research moves model
forecasts toward those contemporaneous prices. It extends an earlier six-contract
model comparison with additional questions; results below concern only the new
cohort.

The experiment crosses three models with two information conditions. In the
\emph{question-only} condition, the worker receives the event definition and
forecasting instructions. In the \emph{outside-research} condition, it also receives
a frozen evidence packet. Calls are separate: the research worker does not see
the question-only response. Thus, a line connecting the two estimates represents
a comparison between calls rather than a conversational revision. There is one
call per model--question--condition cell. Models share the same prompts and
source material within each condition, but retain provider-specific reasoning
settings. The design compares the use of common evidence; it does not compare
models' autonomous search strategies or isolate the effect of every stage in
the broader forecasting package.

\subsection{Selection and chronology}
The candidate inventory contained 20 nominated contracts. Before collecting
quotes, the protocol fixed that inventory, its order, a no-replacement rule,
and eligibility requirements: an active, nonextreme, uncrossed book with spread
at most 10 percentage points and at least one contract available on each side;
the outcome must remain unknown, with at least 12 hours before the relevant
future event. Three candidates lacked a usable two-sided book: September SpaceX
launches, Sam Altman for the Nobel Peace Prize, and \emph{I Love Boosters} for
Best Picture.

Two more were excluded before inference. The current annual SpaceX launch total
could not be established from the accessible primary material, so an already
crossed threshold could not be ruled out. The earthquake definition did not
unambiguously state the start of the counting window. The final cohort comprised
eight macroeconomic contracts, one major-hurricane count, and six NFL games.
No excluded candidate was replaced.

Opening targets were captured around 01:01--01:02 UTC on September 19
(September 18 in New York). Research followed, and the 90 trials were registered
at 01:06:45 UTC. All batches were sealed before target revelation at approximately
01:24 UTC. The largest absolute change between opening and follow-up midpoints
was 0.5 percentage points. This ordering makes the opening midpoint the fixed
primary target; later quotes provide a timing sensitivity rather than a
replacement chosen after observing performance.

\subsection{Evidence and blinding}
Research used primary material from BEA, BLS, the Federal Reserve, the Atlanta
Fed, the National Hurricane Center, and NFL.com. Macroeconomic packets contained
recent releases and relevant projections, with distinctions between historical
measurements, current-quarter nowcasts, and annual projections. NFL packets
contained early-season standings and scoring, scheduled games, and injury
reports. They were deliberately bounded packets, not comprehensive estimates
of team strength. Hurricane evidence relied on captured search-index material
from primary NHC pages because direct access returned HTTP 403; this limitation
was supplied to the workers.

Kalshi provided event definitions and separately stored evaluator targets. It
provided no research evidence or numerical prices to forecasting workers.
Workers had no retrieval tools. For each of the 90 responses, the returned
system and user prompt receipts, scenario identity, and model identity were
checked against the registered inputs. Source packets, jobs, raw results,
accepted forecasts, and evaluator targets were preserved with content hashes.
The coordinator followed procedural separation of the price store; this was
not independent operating-system isolation of that coordinator.

\section{Measures and matched comparisons}
Let $p_{imc}$ denote the probability from model $m$ for question $i$ under
condition $c$, and let $q_i$ be the committed opening midpoint. We report
\[
 \operatorname{MAE}_{mc}=\frac{100}{|S|}\sum_{i\in S}|p_{imc}-q_i|,
 \qquad
 \operatorname{RMSE}_{mc}=100\sqrt{\frac{1}{|S|}\sum_{i\in S}(p_{imc}-q_i)^2}.
\]
Both are in percentage points. For a given pair, the absolute improvement is
$100\{|p_{imQ}-q_i|-|p_{imR}-q_i|\}$. Its percentage improvement divides this
quantity by the question-only absolute gap. A zero initial gap has undefined
relative improvement; negative improvement denotes deterioration. Aggregate
relative improvements below use the ratio of aggregate MAEs, rather than an
average of question-level ratios.

Comparisons use common question sets. The primary research-model comparison
requires an accepted research forecast from all three models ($n=13$). The
paired research-effect comparison requires all six accepted forecasts ($n=9$).
Available-case scores are also retained in the machine-readable results, but
are not used to rank models across unequal question sets. An equal-weight
average of the three research probabilities is evaluated on the same 13
questions as its members.

\begin{table}[htbp]
\centering
\caption{Agreement with the opening midpoint: common research cohort ($n=13$).}
\begin{tabular}{lrr}
\toprule
Model or benchmark & MAE (pp) & RMSE (pp)\\
\midrule
Claude Fable 5.1 & 8.92 & 11.28\\
GPT-6 Astra & 11.08 & 13.84\\
Gemini 3.1 Pro & 13.08 & 15.64\\
Equal-weight research average & 10.87 & 12.68\\
Constant 50\% & 19.38 & 22.41\\
\bottomrule
\end{tabular}
\end{table}

\begin{table}[htbp]
\centering
\caption{Research comparison: fully matched cohort ($n=9$). Distances and absolute reductions are in percentage points.}
\begin{tabular}{lrrrr}
\toprule
Model & Question only & Research & Reduction & Relative\\
\midrule
GPT-6 Astra & 14.56 & 8.00 & 6.56 & 45.0\%\\
Claude Fable 5.1 & 17.33 & 9.56 & 7.78 & 44.9\%\\
Gemini 3.1 Pro & 15.33 & 9.44 & 5.89 & 38.4\%\\
\bottomrule
\end{tabular}
\end{table}

Research improves aggregate agreement for each model, but individual questions
show both gains and losses. For example, all three research forecasts for Q4 GDP
are 42\%, below the 55\% midpoint and farther from it than their question-only
forecasts. All three research forecasts for the December Fed decision are 55\%,
compared with a 30.5\% midpoint. Averaging cannot remove an error shared by all
members. Conversely, research moves the hurricane forecasts close to the 5\%
midpoint. The appended question-faceted slope graphs display each accepted
estimate and its distance reduction, with a labeled horizontal reference at
the midpoint.

The overall model ranking depends on the question mix. Removing NFL contracts
from the common research cohort leaves eight questions: Astra's MAE is 8.44 pp,
Gemini's 9.81 pp, and Claude's 10.44 pp. NFL contracts additionally have a payoff
distinction: a tie pays half, whereas the requested forecast is the probability
of a win. Under a simple expected-payout interpretation, the price corresponds
to $P(\mathrm{win})+\tfrac12P(\mathrm{tie})$, not exactly $P(\mathrm{win})$.
The non-NFL analysis was specified before inference.

\section{Response screening and sensitivity analysis}
All 90 calls completed without provider-reported exceptions. The registered
acceptance rule nevertheless excluded ten responses, leaving 25/30 for Astra,
25/30 for Claude, and 30/30 for Gemini. All ten contain valid JSON probabilities
and valid supplied evidence references. Nine merely deny using market or
sportsbook odds. The remaining response mentions the Kalshi naming convention
when inferring a home team. None of these ten reports numerical market odds;
the lexical matches do not establish price exposure.

The registered exclusions remain unchanged. To expose their effect, a separate
\emph{post-hoc} diagnostic includes all 90 original probabilities, without
editing text, issuing the excluded forecasts, or obtaining replacement calls.
The content review occurred after price revelation and is therefore not part
of the preregistered primary analysis.

\begin{table}[htbp]
\centering
\caption{Post-hoc diagnostic using all original responses ($n=15$). MAE in percentage points.}
\begin{tabular}{lrr}
\toprule
Model & Question only & Research\\
\midrule
GPT-6 Astra & 15.67 & 11.73\\
Claude Fable 5.1 & 15.20 & 9.47\\
Gemini 3.1 Pro & 17.60 & 13.13\\
\bottomrule
\end{tabular}
\end{table}

The equal-weight research average in this diagnostic scores 11.31 pp. The
aggregate direction and overall research ranking agree with the registered
analysis. The screening problem still matters: a future protocol should define
how to distinguish explicit denials from assertions of price use before any
responses are examined. A vocabulary match is useful as an audit flag but is an
imprecise exclusion criterion.

\section{Interpretation and next experiments}
The study provides additional evidence that the supplied research improves
agreement with contemporaneous market assessments on this selected cohort.
The comparison is useful precisely because the events remain unresolved:
realized outcomes cannot explain the observed agreement. Market disagreement
also identifies concrete questions on which evidence selection or numerical
interpretation deserves further examination.

The scope of inference remains the observed questions and settings. A two-sided
quote screen establishes a usable benchmark, not that every contract has deep
liquidity. Related economic releases share information and risks; six games
also do not supply six independent replications of a general forecasting
method. With one draw per cell and a small number of broader topic groups,
we report descriptive differences rather than unsupported precision for a
population ranking. These market-distance measures are intermediate outcomes;
subsequent realized outcomes can support a separate proper-scoring evaluation.

No transformation $f(p)$ or ensemble weights were fitted. Training one on these
same observations and reporting the resulting fit would not establish useful
out-of-sample calibration. The next experiment should freeze any proposed
correction before testing on later, untouched event groups, preserve a simple
identity correction and equal-weight ensemble as controls, and expand coverage
across dates and domains. Research packets and response-screening rules should
be versioned as experimental factors rather than changed within a completed
comparison.

\appendix
\section{Package implementation and reproducibility}
Vorhersage supplies the deterministic records and evaluation machinery around
external research and model judgment. The market workbench stores a question
version and contract definition, checks quote eligibility, and places the opening
quote in a separate evaluator store. A public hash commits to that target before
the research begins. Evidence packets preserve claims, values, source excerpts,
retrieval timestamps, captured content, and declared limitations.

First-class experimental arms combine a method, a model configuration, and a
data specification. Here, six arms cross the three models with either no
current evidence or the same frozen packet for each question. Registration pins
the question versions, information cutoff, arm identities, and one repetition.
The transport exports the exact tasks to independent provider calls. The
importer checks response identity, JSON structure, references, and the declared
acceptance policy before issuing immutable forecast artifacts. Every attempt,
including an excluded response, remains in the usage record.

The reveal gate requires terminal records for every registered trial and seals
for all model batches. It verifies raw-result and forecast commitments before
opening prices. The published comparison then scores all accepted model/condition
forecasts against each question's committed target. A representative forecast
in the workbench journal is only the bridge to that reveal; it is not the source
of the multi-model comparison. An offline verifier independently recomputes
MAE, RMSE, and distance outside the bid--ask spread, and checks target hashes,
retrieval-before-registration timing, and seal-before-reveal ordering.

\begin{table}[htbp]
\centering
\caption{Reported provider inference cost, including excluded responses.}
\begin{tabular}{lrr}
\toprule
Model & Calls & USD\\
\midrule
GPT-6 Astra & 30 & 0.99960\\
Claude Fable 5.1 & 30 & 2.25114\\
Gemini 3.1 Pro & 30 & 0.39297\\
\midrule
Total & 90 & 3.64371\\
\bottomrule
\end{tabular}
\end{table}

Exact identifiers are \texttt{gpt-6-astra}, \texttt{claude-fable-5-1}, and
\texttt{gemini-3.1-pro-preview}. Astra uses high reasoning effort; Claude uses
adaptive thinking with high output effort; Gemini uses a 2,048-token thinking
budget and temperature 0.5. All have an 8,192-token configured output limit.
Serialized legacy sampling defaults may be omitted or ignored by adapters for
reasoning models; the archived configurations and raw provider responses are
the operational record. These settings are not equal compute budgets. The
reported costs exclude research/coordinator work, for which cost is unknown.

Artifacts reside in \texttt{examples/kalshi\_expansion\_20260918/}.
The accompanying README maps protocols, evidence, tasks, jobs, results, seals,
and target disclosures. \texttt{run/comparison.json} contains all cohort
metrics; \texttt{run/output-review-sensitivity.json} contains the post-hoc
review and diagnostic. The saved database integrity check passed for 415
artifacts. The targeted study tests pass 30 tests. No core package dependency
was added for this expansion.

The following commands verify the saved study and regenerate its figures and PDF
from the repository root (with the study dependencies installed):
\begin{quote}\small\ttfamily
.venv/bin/python examples/kalshi\_expansion\_20260918/verify\_results.py\\
Rscript examples/kalshi\_expansion\_20260918/plot\_performance.R\\
latexmk -cd -pdf examples/kalshi\_expansion\_20260918/report.tex
\end{quote}
The offline check needs neither inference calls nor private-store access.
The following ggplot2 plates show the primary comparison: blue circles denote
question-only forecasts, green squares research forecasts, and dashed lines the
opening midpoint. Missing responses are never assigned zero probability.
\clearpage
\includepdf[pages=-,fitpaper=true,pagecommand={}]{figures/performance.pdf}
\end{document}
